\section{Strategi Membangun Bukti Formal Berantai}

Kemampuan inferensi diuji ketika harus menyusun deretan langkah deduktif dari banyak premis menuju konklusi yang ditentukan \cite{rosen2019}. Prosedur tersebut disebut \textbf{bukti formal berantai} (\textit{formal proof of validity}).

\subsection{Prosedur Penyusunan Bukti}

Setiap langkah dalam bukti harus memiliki justifikasi: nomor baris premis atau langkah sebelumnya yang dipakai, serta nama aturan inferensi yang diterapkan.

\begin{contoh}[Argumen majemuk]
Diberikan tiga premis berikut:
\begin{enumerate}
    \item Jika mahasiswa tidak belajar dengan sungguh-sungguh ($\neg p$), maka ia gagal pada kuis ($q$).
    \item Mahasiswa tersebut tidak gagal pada kuis ($\neg q$).
    \item Jika mahasiswa tidak gagal pada kuis ($\neg q$), maka ia berpredikat cerdas ($r$).
\end{enumerate}
\textbf{Tujuan:} Tunjukkan bahwa konklusi $r \land p$ dapat diperoleh secara valid.
\end{contoh}

\subsection{Langkah Pembuktian}

\textbf{Langkah 1: Formalisasi simbolik}
\begin{itemize}
    \item 1. $\neg p \rightarrow q$ \quad (Premis 1)
    \item 2. $\neg q$ \quad (Premis 2)
    \item 3. $\neg q \rightarrow r$ \quad (Premis 3)
    \item \textbf{Konklusi yang dibuktikan:} $r \land p$
\end{itemize}

\textbf{Langkah 2: Bukti berbaris}

\begin{table}[H]
\centering
\small
\begin{tabular}{|p{1.0cm}|p{3.1cm}|p{6.8cm}|}
\hline
\textbf{Baris} & \textbf{Ekspresi} & \textbf{Justifikasi} \\
\hline
1. & $\neg p \rightarrow q$ & Premis 1 \\
2. & $\neg q$                 & Premis 2 \\
3. & $\neg q \rightarrow r$   & Premis 3 \\
\hline
4. & $\neg (\neg p)$          & Modus Tollens pada baris 1 dan 2 \\
5. & $p$                      & Hukum negasi ganda (involusi) pada baris 4 \\
6. & $r$                      & Modus Ponens pada baris 3 dan 2 \\
\hline
7. & \textbf{$\therefore r \land p$} & Konjungsi pada baris 5 dan 6 \\
\hline
\end{tabular}
\end{table}

\textbf{Kesimpulan:} Konklusi $r \land p$ terbukti mengikuti secara valid dari ketiga premis.

Metode pembuktian baris demi baris dengan justifikasi ini disebut \textbf{deduksi formal}. Pendekatan serupa dipakai dalam verifikasi perangkat lunak dan model checking \cite{wiki_first_order_logic}.
